\clearpage\fancyhead[L]{\small KOMET.}\hypertarget{native-394-lower}{}\label{native-394-lower}
\begin{center}\textbf{KOMET.}\par\vspace{4pt}\hrule\vspace{4pt}\fontsize{10}{12}\selectfont Göttingische gelehrte Anzeigen. Stück 149. Seite 1473..1476. 18. September 1815.\par\vspace{4pt}\hrule\vspace{4pt}\end{center}

\par\indent Prof. \textsc{Gauss} hat den höchst bemerkenswerten Kometen, über den wir in dieser Zeitschrift mehrmals berichtet haben, an der hiesigen Sternwarte länger beobachtet, als anfangs zu hoffen war. Da die Beobachtungen jetzt abgeschlossen sind, stellen wir sie hier vollständig zusammen.

\SourceLeaf{395}\begin{center}\fontsize{10}{12}\selectfont\setlength{\tabcolsep}{5pt}\renewcommand{\arraystretch}{1.15}\begin{tabular}{l|c|c|c}
1815 & \shortstack{Mittlere Zeit\\in Göttingen} & \shortstack{Scheinbare\\Rektaszension} & \shortstack{Scheinbare\\Deklination}\\\hline
März 20 & $10^{\mathrm u}\;33^{\mathrm m}\;6^{\mathrm s}$ & $54^\circ\;7'\;1''$ & $39^\circ\;7'\;47''\;\mathrm{N.}$\\
21 & $10\;11\;37$ & $54\;33\;21$ & $39\;36\;57$\\
25 & $9\;46\;38$ & $56\;28\;50$ & $41\;38\;5$\\
30 & $9\;50\;57$ & $59\;13\;3$ & $44\;10\;27$\\
April 2 & $9\;11\;55$ & $61\;2\;27$ & $45\;39\;57$\\
Juni 12 & $10\;53\;31$ & $165\;3\;3$ & $52\;8\;25$\\
30 & $10\;49\;12$ & $184\;39\;24$ & $39\;22\;5$\\
Juli 13 & $10\;57\;20$ & $194\;32\;52$ & $29\;53\;54$\\
27 & $10\;4\;59$ & $203\;0\;56$ & $20\;37\quad:$\\
29 & $10\;13\;58$ & $204\;6\;3$ & $19\;24\;13$\\
Aug. 4 & $10\;14\;7$ & $207\;14\;8$ & $15\;53\;49$\\
25 & $9\;1\;44$ & $217\;1\;31$ & $5\;33\;36$\\
\end{tabular}\end{center}

\par\indent Wir bemerken dazu noch, dass am 13. und 27. Juli und am 4. August Sterne aus \textsc{Piazzis} Katalog zum Vergleich herangezogen wurden. An allen übrigen Tagen wurden dagegen Sterne verwendet, deren Positionen aus der \emph{Histoire Céleste} oder den \emph{Mémoires de l’Ac. de Paris} 1790 reduziert worden waren. Alle Beobachtungen wurden mit dem Kreismikrometer des 10 Fuß langen \textsc{Herschel}schen Teleskops angestellt, ausgenommen die vom 12. Juni. An diesem Tag wurden der Abstand vom Stern Nr. 219 im Großen Bären nach \textsc{Bode} und der Positionswinkel mit einem ausgezeichneten \textsc{Frauenhofer}schen Heliometer gemessen.

\par\indent Die elliptische Bahn, die Herr Prof. \textsc{Gauss} um die Mitte des Juni bestimmt hatte und die im 105. Stück dieser Zeitschrift [S. 385 d. B.] mitgeteilt wurde, stimmte auch mit allen späteren Beobachtungen weiterhin gut überein. Die späteren Bestimmungen dieser Bahn durch die Herren \textsc{Bessel} und \textsc{Nicolai}, die auf einer sorgfältigeren Untersuchung der Beobachtungen beruhen, weichen ebenfalls nur wenig davon ab. Bei der kurzen Umlaufzeit dieses Kometen kann jetzt nur noch eine geringe Unsicherheit bestehen. Wir stellen diese Bestimmungen hier zusammen.

\Needspace{21\baselineskip}\begin{center}
\emph{Elliptische Elemente des Kometen berechnet vom Hrn. Prof.} \textsc{Bessel.}
\end{center}

\begin{center}\fontsize{10}{12}\selectfont Durchgang durch das Perihel am 26. April.\par $0^{\mathrm u}\;5^{\mathrm m}\;14^{\mathrm s}5$ Mittlere Zeit in Paris\end{center}
\begin{center}\fontsize{9}{11}\selectfont\setlength{\tabcolsep}{2pt}\renewcommand{\arraystretch}{1.15}\begin{tabular}{lcr}
Länge des Knotens & $.\;.\;.\;.\;.\;.\;.\;.\;.\;.\;.\;.\;.$ & $83^\circ\;28'\;46''18$\\
Länge des Perihels & $.\;.\;.\;.\;.\;.\;.\;.\;.\;.\;.\;.$ & $149\;2\;29.13$\\
\multicolumn{3}{c}{\shortstack{beide für den 1. Januar 1815 und\\siderisch ruhend vorausgesetzt}}\\
Neigung der Bahn & $.\;.\;.\;.\;.\;.\;.\;.\;.\;.\;.\;.\;.$ & $44\;29\;53.71$\\
Logarithmus des kleinsten Abstands & $.\;.\;.\;.\;.\;.\;.$ & $0.0837950$\\
Exzentrizität & $.\;.\;.\;.\;.\;.\;.\;.\;.\;.\;.\;.\;.\;.\;.$ & $0.93112771$\\
Große Halbachse & $.\;.\;.\;.\;.\;.\;.\;.\;.\;.\;.\;.\;.\;.$ & $17.60964$\\
Siderische Umlaufszeit & $.\;.\;.\;.\;.\;.\;.\;.\;.\;.\;.\;.$ & $73.89682\;\text{Jahre}$\\
\end{tabular}\end{center}

\Needspace{21\baselineskip}\begin{center}
\emph{Elliptische Elemente des Kometen, berechnet von Hrn.} \textsc{Nicolai}, \emph{auf der Seeberger Sternwarte.}
\end{center}

\begin{center}\fontsize{10}{12}\selectfont Durchgang durch das Perihel am 26. April.\par $0^{\mathrm u}\;55^{\mathrm m}\;32^{\mathrm s}$ M. Z. der Seeberger Sternwarte\end{center}
\begin{center}\fontsize{9}{11}\selectfont\setlength{\tabcolsep}{2pt}\renewcommand{\arraystretch}{1.15}\begin{tabular}{lcr}
Länge des Perihels & $.\;.\;.\;.\;.\;.\;.\;.\;.\;.\;.\;.$ & $149^\circ\;3'\;25''3$\\
Länge des Knotens & $.\;.\;.\;.\;.\;.\;.\;.\;.\;.\;.\;.\;.\;.$ & $83\;28\;52.3$\\
\multicolumn{3}{c}{beide für den 26. April und siderisch ruhend}\\
Neigung der Bahn & $.\;.\;.\;.\;.\;.\;.\;.\;.\;.\;.\;.\;.$ & $44\;29\;46.0$\\
Logarithmus des kleinsten Abstands & $.\;.\;.\;.\;.\;.\;.$ & $0.0837490$\\
Exzentrizität & $.\;.\;.\;.\;.\;.\;.\;.\;.\;.\;.\;.\;.\;.\;.\;.$ & $0.93029341$\\
Große Halbachse & $.\;.\;.\;.\;.\;.\;.\;.\;.\;.\;.\;.\;.\;.$ & $17.39704$\\
Siderische Umlaufszeit & $.\;.\;.\;.\;.\;.\;.\;.\;.\;.\;.\;.$ & $72.564\;\text{Julian. Jahre.}$\\
\end{tabular}\end{center}

\Needspace{13\baselineskip}\SourceLeaf{396}\begingroup\setlength{\rightskip}{0pt plus 1em}\par\indent Der Vergleich der letzteren Elemente mit den jüngsten hiesigen Beobachtungen ergab folgende\par\endgroup

\begin{center}
\fontsize{10}{12}\selectfont\setlength{\tabcolsep}{9pt}\renewcommand{\arraystretch}{1.15}\begin{tabular}{l | c | c}
& \multicolumn{2}{c}{Unterschiede} \\
& Rektaszension & Deklination \\
\hline
Juni 30 & $-50''4$ & $+15''9$ \\
Juli 13 & $-16.6$ & $-44.1$ \\
27 & $-23.8$ & \\
29 & $+3.6$ & $-47.8$ \\
August 4 & $-6.5$ & $-38.1$ \\
25 & $-2.4$ & $-47.5$ \\
\end{tabular}
\end{center}

\Needspace{13\baselineskip}\par\indent Bei dieser Gelegenheit holen wir auch den Bericht über die Beobachtungen der \emph{Juno} nach, die Herr Prof. \textsc{Gauss} um die Zeit der Opposition gemacht hat. Dieser diesmal ungewöhnlich lichtschwache Planet scheint an keiner anderen Sternwarte beobachtet worden zu sein.

\begin{center}\fontsize{10}{12}\selectfont\setlength{\tabcolsep}{5pt}\renewcommand{\arraystretch}{1.15}\begin{tabular}{l|c|c|c}
1815 & \shortstack{Mittlere Zeit\\in Göttingen} & \shortstack{Scheinbare\\Rektaszension} & \shortstack{Scheinbare\\Deklination}\\\hline
März 1 & $10^{\mathrm u}\;22^{\mathrm m}\;0''5$ & $197^\circ\;7'\;56''7$ & $2^\circ\;19'\;3''0\;\mathrm{S.}$\\
29 & $10\;8\;57.0$ & $192\;31\;51.8$ & $1\;32\;30.1\;\mathrm{N.}$\\
April 8 & $9\;27\;42.9$ & $190\;34\;51.9$ & $2\;53\;24.2\;\mathrm{N.}$\\
\end{tabular}\end{center}

\Needspace{14\baselineskip}\par\indent Herr \textsc{Nicolai} übernahm es, diese Beobachtungen mit den letzten von Herrn \textsc{Möbius} berechneten Elementen zu vergleichen (siehe Mon. Corresp. 1813, Bd. XXVIII, S. 577). Nachdem er die Epoche um $4'$ $55''$ vergrößert hatte, fand er folgende

\begin{center}
\fontsize{10}{12}\selectfont\setlength{\tabcolsep}{6pt}\renewcommand{\arraystretch}{1.15}\begin{tabular}{l | c | c | c | c}
& \multicolumn{4}{c}{Unterschiede} \\
& Rektaszension & Deklination & Länge & Breite \\
\hline
März 1 & $-20''6$ & $-49.7$ & $0$ & $-53''7$ \\
29 & $-12.2$ & $-50.5$ & $+8.6$ & $-51.2$ \\
April 8 & $-16.3$ & $-45.1$ & $+2.9$ & $-47.8$ \\
\end{tabular}
\end{center}

\noindent und daraus die Opposition mit der Sonne

\begin{center}31. März 1815.\quad $12^{\mathrm u}\;56^{\mathrm m}\;43^{\mathrm s}$ Mittlere Zeit in Göttingen\end{center}
\begin{center}\fontsize{9}{11}\selectfont\setlength{\tabcolsep}{2pt}\renewcommand{\arraystretch}{1.15}\begin{tabular}{lcr}
Wahre Länge & $.\;.\;.\;.\;.\;.\;.\;.\;.\;.\;.\;.\;.$ & $190^\circ\;25'\;6''6$\\
Nördliche geozentrische Breite & $.\;.\;.\;.\;.\;.\;.$ & $6\;29\;10.0$\\
\end{tabular}\end{center}

\Needspace{21\baselineskip}\par\indent Durch die Verbindung dieser Opposition mit den drei vorhergehenden erhielt dieser geschickte Astronom außerdem folgende

\Needspace{18\baselineskip}\begin{center}
\emph{Neue Elemente der Juno.}
\end{center}

\begin{center}\fontsize{9}{11}\selectfont\setlength{\tabcolsep}{2pt}\renewcommand{\arraystretch}{1.15}\begin{tabular}{lcr}
Epoche der mittleren Länge 1816 & $.\;.$ & $230^\circ\;11'\;34''2$\\
\multicolumn{3}{c}{für den Meridian von Göttingen}\\
Tägliche tropische Bewegung & $.\;.\;.$ & $812''9304$\\
Länge des Perihels 1816 & $.\;.\;.\;.\;.$ & $53^\circ\;14'\;53''8$\\
Länge des aufsteigenden Knotens 1816 &  & $171\;9\;58.9$\\
Neigung der Bahn & $.\;.\;.\;.\;.\;.\;.$ & $13\;14\;0.1$\\
Exzentrizitätswinkel & $.\;.\;.\;.\;.\;.$ & $14\;43\;28.84$\\
Logarithmus der großen Halbachse & $.$ & $0.4266844$\\
\end{tabular}\end{center}

\par\indent Eine von Herrn \textsc{Nicolai} aus diesen Elementen berechnete Ephemeride wird im Astronomischen Jahrbuch für 1818 veröffentlicht werden.
